\chapter{图像和多模态处理概述}
\label{chap:vision-overview}

同一段街道录像可以支持许多不同的请求。交通统计关心经过多少辆车，图像检索希望找回拍到某辆车的其他照片，夜景增强希望看清路边招牌，视频编辑则可能要求把车身改成另一种颜色。这些任务都读取视觉信息，但需要保留的证据、允许改变的内容和应当返回的结果并不相同。仅仅选择一个能够接收图像的模型，还没有完整规定要解决的问题。

前面的模型与范式介绍了表示、预测、生成和学习的通用机制。本卷从应用端继续追问：观测怎样进入计算，模型输出怎样还原成对象、区域、时间区间或媒体成果，怎样判断结果满足了最初的请求？理解本章只需知道训练利用样本调整预测规则，使用时再把这条规则作用于新输入；具体网络尚不是阅读前提。本章先规定观测、成果和约束，再说明这些要求怎样分配到全卷五个部分十五章：第一部分在本章之后讲图像、视频和多模态数据的经典基础，第二部分的\ref{chap:vision-image-video-models}章和\ref{chap:vision-multimodal-models}章建立任务共用的图像、视频和多模态模型，后三个部分再分别加入图像中的空间要求、视频中的时间要求和跨来源的对应要求。这样的安排使我们先知道需要什么结果，再研究哪些计算能够形成它；研究历史和当前前沿则说明这些要求是怎样逐步进入研究、目前还缺什么。

\section{图像和多模态处理的基本概念}
\label{sec:vision-overview-concepts}

描述一个处理问题，需要说清两件事：系统依据什么，以及系统交出什么并受什么约束。前一项包括观测本身，也包括标注、查询、创作条件等不同职责的信息；后一项包括成果的形式，以及能检查的约束量、坐标与时间的还原方式和评价依据。把两者混在一起，容易把“图像变清晰了”误认为“图中文字恢复正确了”，或把“能输出一个框”误认为“找到了请求中的对象”。

以“统计这段街景中转弯的车辆”为例，完整请求还要说明统计哪段时间、哪个路口，以及同一车辆多次出现怎样计数。系统需要取得相关画面，保留拍摄时刻，把局部预测对应到原画面，再判断哪些对象完成了约定的动作，最后按个体而不是按帧计数。这是一条有输入输出依赖的处理过程：没有时刻便难以确定事件区间，没有个体对应便可能重复计数。具体检测和跟踪算法可以不同，这些任务关系却都需要明确。下面先说明输入一侧有哪些信息、各自承担什么职责，再说明输出一侧怎样规定成果并使之可以核验。

\subsection{观测与输入依据}
\label{sec:vision-overview-observations}

街景照片没有直接存下“汽车”“招牌”这些名称，而是把成像结果记录在离散网格上。\term{像素}（pixel）是网格中的一个位置及其记录值，\term{通道}（channel）则区分该位置记录的不同分量。常见的红绿蓝图像（red, green, blue, RGB）在每个位置记录三个颜色分量，但通道也可以承载深度、光谱或其他测量。为了说明算法接收多少位置、每个位置有多少数值，把图像写为三阶数组
\begin{equation}
  \bm{X}\in\mathbb{R}^{h\times w\times c},
  \qquad x_{i,j,k}=\text{第$i$行、第$j$列、第$k$通道的值。}
  \label{eq:vision-image-array}
\end{equation}
这里$h,w,c$分别是高度、宽度与通道数。数组字形沿用本书的粗体约定，元素值是标量。输入范围也属于定义的一部分：8位整数的$0$至$255$与归一化后的$0$至$1$不能直接混用。

例如，把一幅8位RGB照片交给要求$[0,1]$输入的程序，需要按约定将各分量除以255；这改变了数值尺度，不改变采样位置。若程序要求先减去通道均值，还要记录相应常数。训练和使用必须采用相容的输入约定，否则同一照片会形成不同数值，不能把由接口不一致造成的变化都归因于场景不同。

\term{分辨率}（resolution）在数字图像接口中通常指网格尺寸，例如$1920\times1080$；它决定有多少采样位置，却不能独自说明可辨细节。失焦、高噪声或压缩过重都可能使较大的图像缺少有效信息。镜头、曝光、光照、传感器和后续成像处理影响像素值，同一物体在不同照片中的像素未必相同。把一幅小图插值到更大网格，也没有增加新的实拍观测。

像素坐标只是图上的位置。本卷在没有另行说明时使用横向向右、纵向向下的图像坐标；连续框边界用两个坐标之差计算宽高。若讨论离散像素中心、图像缩放或归一化坐标，会明确给出变换。对象身份用于说明不同位置是否属于同一实体，三维坐标用于说明实体位于哪个空间位置。两辆车可在不同时间经过同一像素，一辆车也可以跨越许多像素，因此这三种描述不能互换。

\term{视频}（video）把画面与时间关联起来，也就是记录不同拍摄时刻的视觉观测。为了区分画面顺序和实际时间，可表示为$\{(\bm{X}_t,\tau_t)\}_{t=1}^{n}$，其中$n$是帧数，帧编号$t$表示顺序，时间戳$\tau_t$以秒为单位表示拍摄或播放时刻。帧率（frame rate）$f$表示每秒记录多少帧；固定帧率时，相邻帧间隔为$1/f$，抽帧步长为$k$时，被保留的相邻画面间隔为$k/f$。这两个间隔分别属于原始记录和实际输入，比较运动或事件时不能混用。

可变帧率或丢帧时，应使用真实时间戳，不能简单把帧编号当秒数。镜头切换会中断局部运动的连续性，相机移动也会让静止物体在图上移动。例如，相机向右平移时，路灯可在画面中向左移动；只凭这个像素位移，无法断定路灯正在运动。视频的空间变化需要结合相机、对象对应和时间条件解释。

文字、图像和声音还可以共同描述街景中的一个事件。\term{模态}（modality）指信息的类型或通道，例如视觉画面、文字表达和声音记录。\term{多模态处理}（multimodal processing）利用不同类型信息之间的联系完成任务，也就是让一种信息帮助选择、补充或核查另一种信息。画面可以显示车辆的运动，声音可以提供鸣笛时刻，文字可以指定应当寻找哪次转弯；这些联系需要学习或建立对应，不能由“同时输入”自动得到。

同一辆车的多张照片增加的是视角和观测次数，图像与文字配对增加的是信息类型及两者之间的联系。视频虽然含有许多帧，在本卷仍按时序视觉输入展开；页面布局说明文字、图像与表格怎样排列，也不必另立为一种模态。这样的约定用于确定章节职责，不宣称各种文献都采用完全相同的模态划分。

\begin{example}[采样之后还能看见什么]
  \label{ex:vision-overview-sampling}
  以下为构造的教学情形。设视频按区间$[0,12)$秒记录，固定帧率为25帧每秒，含300帧，第$t$帧的时间为$(t-1)/25$秒。保留第1、26、51帧并依此每隔25帧取一帧，得到0至11秒处的12幅画面，间隔为1秒。设一辆车只在4.2秒至4.6秒间出现在画面内，则这些输入中完全没有这辆车。后续模型即使正确处理每幅输入，也不能据此证明原视频中没有车辆出现。
\end{example}

这个例子把信息获取和模型计算分开了。错误可能发生在进入模型之前；增加语言推理长度不能补回从未保留的那段观测。长视频、高分辨率页面和细小目标都需要检查同样的问题。

观测之外，处理问题中的其他信息还有不同职责。人工类别、框和掩码是训练或评价标注，查询说明希望找什么，创作条件说明希望生成什么，相机参数与曝光记录则提供辅助测量。一句“红色汽车正在转弯”若是检索查询，不能同时被当作视频中确有这一事件的证据；若是创作条件，它规定的是希望出现的内容，而不是已经发生的事实。

这些信息的取得过程也不同。训练时，一条样本可以包含街景图像、车辆范围和类别标注，算法据此比较预测与目标，再调整模型中可学习的参数。使用时，人工正确框不再作为输入，系统必须从新观测形成候选并按任务规则选择。相机参数来自标定，查询来自用户，模型权重来自训练；这些信息的来源不同，不应都称为“模型自己知道”。具体误差、参数更新及停止条件由模型章和各任务章展开。

配对信息也可能不完整。网页文字常只描述图片的一部分，旁白可能先讲下一步操作，字幕可能来自画外人物。需要把图像编号、区域、时间与来源保留下来，再检查是否对应同一实体或事件。模型得分度量的是给定训练与接口下的匹配程度，不是对配对事实的无条件证明。

因此，一条可回查的结果至少要说明使用了哪份原记录、哪段时间和哪个区域，输入经历了怎样的裁剪或采样。若请求要求“转弯时鸣笛的车”，系统既要找到转弯范围，也要取得对应录音，并核对声音是否能归于这辆车。只有同期鸣笛声而没有声源依据时，可以报告观察到的两件事及其时间关系，不能自动补出同一主体，更不能由先后顺序推出原因。来源记录让输入依据在输出端仍然可以追溯，这一要求在下一小节落实为结果的还原。

\subsection{输出形式与结果约束}
\label{sec:vision-overview-outputs}
\label{sec:vision-overview-evidence}

输入依据确定之后，还要规定系统交出什么。当请求是“图中有没有公交车”时，结果可以是类别判断；请求变为“公交车在哪里”时，还需要空间范围。\term{目标检测}（object detection）通常返回对象的边界框（bounding box）、类别与得分，边界框用矩形圈出对象所在范围。\term{图像分割}（image segmentation）则为像素指定类别或对象归属，也就是在许多位置分别说明那里是什么或属于谁。框包含的背景多少、两个同类对象是否分开，都影响结果含义。

分割结果可以用掩码（mask）表示，即逐位置记录哪些像素属于目标。例如，一个框可以圈住整辆公交车并包含周围背景，车辆掩码则可以沿轮廓排除背景。若只给每个像素“车辆”类别，两辆相邻车辆仍可能连成同一片区域；若要求分别编辑两辆车，就还需区分它们的个体归属。对象跟踪（object tracking）进一步把不同时刻的范围关联为同一对象的轨迹（trajectory），即同一对象在不同时刻的位置序列，不能用一串没有身份对应的框替代。

数量、关节位置和距离属于测量结果。按检测实例计数需要规定怎样处理边界上的对象；从图像估计距离则需要交代相机、尺度和坐标依据。检索返回满足相关性要求的记录及排序。“外观像同类汽车”和“拍到了同一辆汽车”是不同的相关性，不能用同一组正例默认替代。

空间测量还要区分要测哪个量。姿态（pose）可以描述对象的关节位置或朝向，深度（depth）描述沿约定相机方向的远近；两者既不同于类别，也不同于框的宽度。若没有已知尺度或其他标定依据，单幅照片中的“大车”和“近车”可能给出相近的像素尺寸。语言描述“车离路灯很近”也没有直接给出以米计的距离，测量要求不能由自然语言措辞替代。

另一类请求需要返回图像本身。\term{图像复原}（image restoration）依据退化观测估计原内容，例如从带噪照片估计较少噪声的图像；\term{图像增强}（image enhancement）改善指定用途下的呈现，例如调整暗部层次；\term{图像压缩}（image compression）用较少比特表示图像，并规定解码后的允许误差。三者可以组成一条处理链，但分别回答内容估计、呈现和表示成本的问题。

压缩成果首先是码流（bitstream），也就是按规定组织的一串比特；接收端通过解码把它还原为图像。只比较解码图像是否清晰，遗漏了码流大小；只比较文件是否小，遗漏了失真及解码条件。增强成果则需要按用途核验，例如招牌能否读清、暗部层次是否便于观看；这两项要求可能选择不同的亮度变化。

\term{图像生成}（image generation）根据分布或条件创建图像，\term{图像编辑}（image editing）则对已有图像实施指定变化。复原与生成都可能产生原输入中看不到的细节，区别在于任务要求：复原希望估计原场景，创作可以按要求改变场景。一个去模糊模型把招牌上的字改成另一个清晰的字，即使观看效果很好，也没有完成文字保真的复原目标。

理解和生成因此不是按“输出文字还是输出图片”划分。回答“车身是什么颜色”虽然生成了文字，任务仍是依据观测作出判断；按请求把车身改成蓝色，任务则允许改变原内容。共享一个模型并不取消这一区别，因为两个请求允许使用的依据和允许发生的变化不同。

图\ref{fig:vision-task-map}把同一街景可能形成的成果并列起来。请沿各行比较结果需要对应的语义、范围、身份、尺度或内容条件，而不是把各行理解为互斥的模型类别。共享网络可以接不同输出模块，一个模块预测框，另一个模块组织描述；图中也把描述与编辑、配声放在同一成果分支，但描述仍需观测支持，不因此具有任意创作的权限。

\begin{figure}[htbp]
  \centering
  % 教学任务地图；宽约108mm，采用项目语义配色。
\begin{tikzpicture}[
  font=\figurefont\small,
  task/.style={draw=BookRule,fill=BookPaper,minimum width=30mm,
    minimum height=9mm,align=center,inner sep=2mm},
  link/.style={-{Stealth[length=2mm]},draw=BookMuted,line width=0.8pt},
  note/.style={anchor=west,text=BookMuted,font=\figurefont\footnotesize}
]
  \node[task,fill=FigureInput!10,draw=FigureInput,minimum width=20mm] (input) at (0,0) {街景观测};
  \node[task] (recognition) at (4,2.6) {类别、属性};
  \node[task] (region) at (4,1.3) {框、掩码、轨迹};
  \node[task] (measure) at (4,0) {数量、距离、排序};
  \node[task] (restore) at (4,-1.3) {复原、压缩结果};
  \node[task] (create) at (4,-2.6) {编辑、描述、配声};
  \draw[draw=BookMuted,line width=0.8pt] (input.east) -- (1.8,0);
  \draw[draw=BookMuted,line width=0.8pt] (1.8,-2.6) -- (1.8,2.6);
  \foreach \name/\y in {recognition/2.6,region/1.3,measure/0,restore/-1.3,create/-2.6}{
    \draw[link] (1.8,\y) -- (\name.west);
  }
  \node[note] at (5.8,2.6) {对象与语义};
  \node[note] at (5.8,1.3) {范围与身份};
  \node[note] at (5.8,0) {尺度与相关性};
  \node[note] at (5.8,-1.3) {保真与比特};
  \node[note] at (5.8,-2.6) {条件与对应};
\end{tikzpicture}

  \caption{同一街景观测可以服务不同任务。各行表示一种成果及其主要核验关系，均为教学构造，不代表某个模型的实测输出。图像、视频、文字和声音可以分别进入相应分支；输入相同不意味着监督或评价相同。}
  \label{fig:vision-task-map}
\end{figure}

成果形式确定后，约束必须落实到可以检查的量。车辆计数可以规定在某时间段越过指定线的个体数；文字读取可以要求完整字符串一致；压缩可以给定每像素比特数；交互分割可以限定用户点击次数。分辨率、允许等待的时间和候选选择预算会改变可用信息，也必须进入比较协议。

以12秒的街景计数为教学例子，若任务只统计在$[0,12)$秒内越过路口入口线的车辆，那么停在路边却未越线的车辆不计入，同一车辆连续出现在20帧中也只计一次。边界时刻、遮挡时的处理及无法判定身份时的输出都应预先约定。这样得到的数值才有确定的统计范围和对象；“检测到20个框”还不能直接当作“经过20辆车”。

形成最终成果也要撤销输入组织造成的坐标与时间变化，并保留不能恢复的部分。模型看到的网格可能只是原图的一部分：若先裁出招牌再缩小，预测的文字区域就处于裁剪图坐标中，返回原图位置时需要撤销缩放并加回裁剪起点；只保存裁剪后的图片，会丢失它原来位于哪幅照片、哪个位置的依据。同样，抽帧得到的事件边界需要按时间戳转换，码流需要由规定解码器还原，语言答案需要回指支持它的记录。由二维位置换到三维位置则是另一类问题，需要投影关系、相机条件及尺度信息，不是把横纵坐标再添一个数就能完成。若采样遗漏了关键时刻，或解码约定不匹配，应说明证据不足或接口错误，不能用一个看似完整的输出掩盖缺口。

评价依据最后决定一个结果能代表多大范围。数据划分应对应实际使用中的独立单位。如果相邻视频帧随机分入训练与测试，两边可能含有几乎相同的背景和对象，测得的结果就不能充分代表面对新视频的表现。医学图像常需按患者或采集中心划分，工业图像可能需要按批次与设备划分，文档可能需要按来源与模板划分。选用哪一种划分取决于希望估计的未来变化。

若部署目标是另一地区的路口，按地区留出数据比混合各地相邻帧更接近这一检验问题；若目标是在同一路口处理新设备，则又需要按设备或采集批次检查。验证数据用于选择设置，最终测试数据用于估计选定系统面对留出条件的结果，两者职责不能混用。这些划分不保证覆盖全部未来情况，但决定当前评价可以支持哪种外推。

ImageNet-C是分类扰动基准，给既有图像施加规定类型和严重度的合成退化；WILDS是分布变化基准集合，把真实数据中的环境差异组织成训练与评价条件，不同数据集采用不同分组。二者提供的证据不同：在合成噪声上类别错误较少，并不能直接推出跨医院或跨地域也有同样结果。这里由像素退化扩展到数据来源变化，原有结论只能保留在已经检验的退化类型内\scite{hendrycks2019corruptions,koh2021wilds}

\section{图像和多模态处理的研究内容}
\label{sec:vision-overview-tasks}

上一节的观测、成果和约束可以组合出很多应用，研究内容需要按某种依据分工。本卷按主要成果及新增信息条件安排五个部分十五章。第一部分在本章之后用五章讲经典基础，它们既是后续学习模型的设计来源，也是任务章中与学习方法对照的路线：\ref{chap:vision-image-formation}章说明图像怎样由光照、传感器和成像管线形成，怎样采样与量化，以及怎样用色彩空间和相机几何描述；\ref{chap:vision-image-transforms}章研究灰度变换、空间滤波、频域处理、多分辨率表示和形态学；\ref{chap:vision-image-features}章研究边缘、局部特征、形状描述、经典分割和基于人工特征的识别；\ref{chap:vision-video-foundations}章研究视频的时间采样、运动估计、结构分析和时间连续性的度量；\ref{chap:vision-multimodal-alignment}章研究多模态数据的特点、跨模态对齐、融合以及模态缺失与不平衡下的学习。第二部分建立后续任务共用的模型：\ref{chap:vision-image-video-models}章“图像与视频模型”说明视觉输入怎样被编码成空间或时空表示，任务输出怎样由全局向量、空间特征图或稠密预测头解码，并介绍图像理解、视频理解和视觉生成的代表模型；\ref{chap:vision-multimodal-models}章“多模态模型”在视觉表示之上加入文字和声音，说明不同模态怎样连接，以及怎样形成匹配分数、语言响应和媒体成果。

模型可以跨任务复用，任务却各自带来新的要求。第三部分处理图像任务，在单幅或多视图观测上形成类别、范围、测量、复原结果和新图像；第四部分处理视频任务，增加跨帧对应、事件时间和信息到达条件；第五部分处理多模态任务，组合查询、证据与多种输出。同一模型可以出现在多个部分，章节归属取决于当前主要问题，不取决于模型名称。下面三个小节分别说明三个应用部分研究什么、面对哪些主要挑战；每章的挑战在该章第一节展开，经典的应对办法来自第一部分的基础章。

\subsection{图像应用的研究内容}
\label{sec:vision-overview-image-tasks}

在工业质检中，“这个零件是否异常”和“异常位于哪里”是两个相互联系的任务。正常样本的局部特征可以给出偏离程度，区域结果还需与具体缺陷对应。若照片本身严重模糊，复原可以改善可见性，但生成出的划痕不能作为真实缺陷的新证据。图像识别、区域判断与复原应在相同对象上保留来源关系。

在这组问题中，识别把局部和整体线索变成类别、属性或异常判断，检测和分割把判断落到范围，姿态和空间恢复再形成关节、相机、深度或三维结构。空间恢复（spatial reconstruction）利用图像等观测估计场景的空间组织，不等于任意创建一个好看的三维场景。\ref{chap:vision-image-recognition}章沿这些成果进入测量及匹配，解释每一步需要怎样的监督、坐标和尺度依据。

照片整理提供另一组联系。场景类别可以先缩小候选，再由整图表示找相似照片，局部对应与几何核验进一步检查是否拍到了同一建筑。按“街道”“夜景”整理是类别检索，按“红色车身”整理是属性检索，找回同一辆车或同一建筑则是实例检索（instance retrieval），即把具体实体而不是同类外观作为相关性依据。最后一种请求需要保留能够辨认个体的细节，不能只靠场景名称。

\ref{chap:vision-image-restoration}章分别研究退化观测中的内容估计、指定用途下的呈现及有限比特表示。\ref{chap:vision-image-generation}章研究按条件创建、输入到输出的内容转换、局部或连续编辑，以及已观测场景的新视角呈现。对于工业零件，恢复原有刻痕与创作另一种表面纹理具有不同目标；对于照片整理，压缩后仍能区分个体与只保留场景大意也具有不同要求。质量、测量和生成因此可以协作，但需要分别说明哪些细节可改变。

图像任务的挑战首先来自观测本身。成像会模糊、加噪、混叠和量化，不同的原场景可能给出同一幅观测，复原因此是病态逆问题，需要额外的先验才能在多个候选中作出选择；同一物体又随视角、光照、尺度、遮挡和形变呈现截然不同的像素，而“汽车”“缺陷”这类语义与像素值之间隔着很长的推理距离，常称为语义鸿沟。另一些挑战来自输出和数据：分割、深度等稠密输出要求逐像素给出结果，高分辨率输入却受计算预算限制，缩小输入又会丢掉小目标；逐像素和逐框的标注昂贵，类别分布呈长尾，罕见缺陷只有少量样本；训练与部署所用的相机、天气、地区不同，还会带来域偏移。

经典方法用成像模型和正则约束处理逆问题（\ref{chap:vision-image-formation}章、\ref{chap:vision-image-transforms}章），用人工设计的不变特征和区域划分应对外观变化（\ref{chap:vision-image-features}章）。学习方法从数据中获得先验和表示，替代了大量人工规则，但没有取消观测的限制：先验补出的细节仍可能不是原场景的内容，训练分布之外的外观仍可能失效。识别、复原和生成三章分别在\ref{sec:vision-recognition-challenges}节、\ref{sec:vision-restoration-challenges}节和\ref{sec:vision-generation-challenges}节说明这些挑战在各自任务中的具体形式。

\subsection{视频应用的研究内容}
\label{sec:vision-overview-video-tasks}

看到每一帧中的车辆还不足以统计车辆经过路口的数量。统计需要区分同一辆车的重复出现，跟踪需要在遮挡前后维持身份，转弯判断又要把轨迹与道路及时间范围联系起来。逐帧检测正确但身份交换，是跟踪错误；身份正确但事件边界偏移，则是活动定位错误。

对同一段车辆转弯录像，轨迹回答“同一辆车在何处”，事件区间回答“何时发生了约定动作”，检索结果回答“哪份录像或哪段片段符合查询”。活动和事件不是仅凭一帧的类别就能确定的量，例如转弯需要观察方向随时间改变，并结合道路关系解释。\ref{chap:vision-video-recognition}章“视频感知、跟踪与检索”据此连接对象和属性、范围延续、活动与事件、视频及片段检索，不把这些成果合并成一个逐帧准确率。

邻帧也可以提供单幅图像缺少的细节。某帧的路牌模糊，另一帧可能清晰；只有把两帧中同一位置的内容对应起来，额外观测才可能帮助恢复。把对应内容放到可比较位置称为对齐（alignment），把有效信息传给其他帧称为传播（propagation）。更多帧同时带来等待、存储与错误传播；若输出时不允许等待后续帧，就不能直接使用双向读取前后帧的算法。

超分辨率（super-resolution）估计更密空间网格上的内容，插帧（frame interpolation）估计已有两帧之间的画面，二者分别扩大空间和时间采样。例如，暗处的车辆若在相邻帧中可见，正确对应可以利用这条补充信息；若误把另一辆车的纹理传过来，单帧虽然更清晰，主体却错了。\ref{chap:vision-video-restoration}章讨论这些跨帧恢复、采样扩展及编码问题，清晰度、身份延续与时间信息分别核验。

视频编码还要规定编码端使用哪些历史、向解码端传递什么，解码端怎样仅凭码流和允许的参考帧恢复画面。离线处理（offline processing）可以读取整段给定录像，在线处理（online processing）按输入到达顺序工作；两种方式都可能维护长历史，但能否读取相对于当前时刻的后续帧不同。“输入帧数更多”“历史更长”和“当前输出更准确”分别描述信息量、保留范围和任务结果，不构成彼此的保证。

视频在图像的全部挑战之上增加了时间维度的困难。逐帧独立处理的结果会闪烁，同一对象的身份会在交会时交换，时间上的连续与一致需要专门维护；画面中的像素位移同时包含对象运动和相机运动，不分开二者就不能解释谁在动；事件的时间尺度从几帧到几分钟不等，长程依赖要求保留远处的历史，而相邻帧高度冗余，逐帧计算和全部保存都超出预算。在线应用还要求因果和低延迟，当前输出不能读取未来帧；动作何时开始、何时结束本身就有歧义，不同标注者给出的时间边界并不一致。

经典方法用块匹配和光流描述运动，用背景建模分离前景，用镜头检测和关键帧压缩冗余（\ref{chap:vision-video-foundations}章）；学习方法用时空表示、跨帧注意力和记忆结构把这些步骤并入网络，但更长的输入并不自动带来更可靠的时间判断。\ref{sec:vision-video-recognition-challenges}节和\ref{sec:vision-video-restoration-challenges}节分别说明这些挑战在视频感知与视频复原中的具体形式。

\subsection{多模态应用的研究内容}
\label{sec:vision-overview-multimodal-tasks}

“找到车辆转弯的片段”把文字查询与时间证据联系起来；“为什么画面中有人回头”可能还需要听见声音，但可观察到先后关系并不等于识别了因果机制。文档和图表增加了字符串、阅读顺序、行列与单位的对应。\ref{chap:vision-multimodal-understanding}章从检索、定位、读取和结构恢复推进到描述与问答，始终保留答案所依赖的来源。

例如，“结合声音说明转弯时发生了什么”需要先对应画面与录音的时间，区分发动机声、鸣笛声和画外声，再形成有来源的描述。文档中同一个数字若属于不同表头，含义也可能不同；界面上的同一个图标若处于不同面板，指向的操作又可能不同。多图、文档、图表、界面和音画因此不是仅增加几个输入文件，还增加了跨页、跨区域、跨时刻的对应要求。结构化信息（structured information）把读到的内容组织为字段及明确关系，例如把数字与其行名、列名和单位一同返回，而不是只输出一串字符。

“生成同一辆车转弯并带发动机声音的片段”则改变了问题性质。这里的车、运动和声音是创作条件，不是待解释的实拍证据。\ref{chap:vision-multimodal-generation}章讨论视频、配声、联合音视频及空间成果的制作。固定已有视频再生成声音，与画面和声音共同生成，具有不同的可修改范围；单幅文生图虽然也是跨模态任务，其完整应用流程集中在图像生成章。

制作多个成果时，还需要说明它们共享什么条件。车辆参考图规定主体外观，文字规定动作，已有录音可以规定节奏，但参考图并没有规定未见一侧的全部结构。视觉条件响应（visually conditioned response）把画面中的判断表达为文字或语音，也就是根据所见内容组织回答，仍需核查事实，并按播放要求组织输出。纯文本处理和不依赖视觉的语音应用属于卷九的范围；本卷只在视觉证据或媒体协调改变问题时展开相应关系。

多模态处理的挑战来自模态之间的差异与对应。像素网格、文字词元和声音波形的结构、采样率和统计性质都不同，表示是异构的；对应关系又有多种粒度，可以是整图与整句、区域与短语，也可以是视频帧与单词、声音事件与画面时刻。网络收集的图文配对往往很弱，文字只描述图像的一部分，甚至与图像无关。融合时，容易学习的模态会主导训练，另一模态的信息被忽视，形成模态不平衡；使用时某一模态还可能缺失，或者与其他模态相互矛盾。生成任务则是一对多的：同一句描述对应无数合理的图像或视频，难以只用一个参考结果评价。

经典方法用典型相关分析、动态时间规整建立对齐，用早期或晚期融合组合各模态（\ref{chap:vision-multimodal-alignment}章）；学习方法用对比学习在大规模弱配对数据上学习共同表示，用交叉注意力在中间层建立细粒度交互，但模型给出的匹配分数仍不是配对事实的证明。\ref{sec:vision-understanding-challenges}节和\ref{sec:vision-multigen-challenges}节分别说明这些挑战在多模态理解与生成中的具体形式。

\section{图像和多模态处理的研究历史}
\label{sec:vision-overview-history}

从局部图像匹配到连续音画响应，应用扩展不断增加需要核验的对象和关系。回顾历史要回答的问题是：观测约束、输出组织、开放条件和连续多源输入这些要求分别在什么时候、以什么形式进入了计算。视觉研究没有一条覆盖所有任务的单线发展史，分类、几何、复原、编码与生成长期并行，也不断借用彼此的表示和计算。以下按应用要求的扩展划分四个阶段，阶段以代表工作的公开年份标注，各路线在阶段边界前后仍然延续；代表工作只用于解释新增的问题。

\subsection{人工设计方法与观测约束（2012年以前）}
\label{sec:vision-overview-classical-history}

同一街景在两次拍摄中会改变大小、方向和亮度，恢复或匹配不能只比较原像素是否相等。早期及长期使用的人工设计路线把所需关系写成明确规则：滤波（filtering）按邻域规则组合像素，成像与退化建模说明场景怎样变成带模糊或噪声的记录，多视图几何（multi-view geometry）利用不同相机投影之间的约束恢复空间关系。已知模糊过程时，可以研究哪些变化被抹弱及怎样限制恢复时的噪声放大；多幅照片有可靠对应时，可以利用投影关系估计空间结构。观测模型也由此指出输入无法唯一确定哪些内容。

尺度不变特征变换（scale-invariant feature transform, SIFT）的2004年论文提供了局部匹配的代表入口：它寻找图像中在不同尺度下可重复辨认的位置，为位置确定参考方向，并把附近外观整理为可比较的描述。这样可以从两幅建筑照片提出局部对应，再由几何约束核查。名称中的尺度不变性指设计所针对的尺度变化，不是任意视角、遮挡或材质变化下的匹配保证\scite{vision-lowe2004}

这类方法没有因深度网络出现而失去全部作用。几何一致性仍可检查学习式匹配结果，退化与采样关系仍可约束复原，码流解码仍须遵守编码端与解码端的信息一致性。后来变化较大的是哪些表示与计算由数据学习，物理观测的限制并未随之消失。

\subsection{学习式视觉应用的扩展（2012—2019年）}
\label{sec:vision-overview-learned-history}

若任务是从许多照片判断类别，人工规定每种外观线索会遇到覆盖和组合的困难。学习式路线让训练样本同时约束内部表示和最终输出。AlexNet是2012年大规模图像分类的代表网络，它把像素变换成多层表示，再给出整图类别分数；训练中的图像类别监督说明应当选择什么类别，却没有直接告诉模型对象的准确范围\scite{krizhevsky2012}

要判断路面、车辆和建筑分别占据哪些位置，整图类别就不够了。全卷积网络（fully convolutional network, FCN）的代表工作于2014年提交预印本、2015年发表于CVPR，把预测组织为与空间位置对应的类别图；这种密集输出（dense output）在许多位置分别给出结果，而不是只返回一个整图标签。它使学习到的图像表示进入像素级任务，同时需要与空间范围相对应的监督\scite{vision-1411-4038}

检测进一步返回多个对象范围，实例分割区分同类个体，全景分割（panoptic segmentation）共同组织可数对象的个体区域与路面、天空等背景区域。这里增加的是结果怎样组织、哪种标注能够核验它的问题。视频还要跨帧关联同一对象，不能由单幅分类发展独自概括这些进展，也不能只用“网络变深了”解释输出职责的变化。同一时期加深网络的残差结构也使分类网络成为许多任务共用的特征提取部分，这种用法在\ref{sec:vision-ivm-resnet}节展开。

与此同时，视频中的时空表示和对象关联、低照度成像、图像转换及学习式压缩也在发展。它们分别处理时间证据、成像条件、输入输出映射和有限比特表示。一个网络能够在多种任务中使用，不代表同一检查点在没有对应监督或条件接口时就具备所有能力。

\subsection{可复用能力与开放条件（2020—2023年）}
\label{sec:vision-overview-open-history}

这一阶段的一项变化是让结果组织和使用条件更加开放。检测Transformer（Detection Transformer, DETR）的2020年工作把检测结果看作无序对象集合：训练时将预测与标注逐一匹配，使集合中的元素对应不同目标。这里的无序指“第一个框”没有固定类别或身份，不指对象之间没有关系；检测的具体训练与输出选择在识别章展开\scite{vision-2005-12872}

对比语言图像预训练（contrastive language-image pre-training, CLIP）的2021年工作让图片和文字形成可比较的表示，用户因此可以在使用时写出候选类别描述，再比较新图像与各描述的分数。开放类别（open-vocabulary categories）在这里表示候选不再只能由固定分类输出表规定，不意味着任意新概念都能识别正确。CLIP没有仅凭这种匹配接口就获得自由句子生成能力\scite{radford2021}

生成式视觉语言模型（generative vision-language model）则把视觉信息接入语言预测，依据图像与请求组织描述或答案。与CLIP选候选不同，它需要逐步形成可变长度文字，也增加了核查每条回答是否由观测支持的要求。两类接口及其训练依据由\ref{chap:vision-multimodal-models}章建立，历史在这里关注它们让用户提出了哪些新的请求。

用户还可以用点或框选择待分割对象，用空间条件约束生成布局，用少量参考图像指定创作主体。提示扩展了交互方式，也引入了目标歧义：落在车门上的一点可以指车门，也可以指整辆车。开放词汇和通用接口扩大了候选范围，结果仍须依照具体目标核验。

新视角合成（novel-view synthesis）形成了另一条支线：给定已拍照片及相机条件，估计从另一位置观看这个场景会得到什么画面。2020年的神经辐射场（neural radiance field, NeRF）用可查询的空间密度和方向相关颜色表示场景，2023年的三维高斯溅射（3D Gaussian Splatting）用一组空间高斯元素表示可渲染内容。两者依据已拍视图拟合场景，再计算新相机下的图像；未观测区域仍可能受先验或表示约束影响，照片呈现准确也不自动证明表面几何适合测量。相关机制在\ref{sec:vision-generation-novel-view}节展开\scite{vision-nerf2020,vision-gaussian2023}

\subsection{连续、多源与交互应用（2024年以来）}
\label{sec:vision-overview-continuous-history}

用户在第一帧指出一辆车后，往往还希望在后续画面中继续获得它的区域。分割一切模型2（Segment Anything Model 2, SAM 2）于2024年公开，这里依据同年10月的v2说明它怎样将点或框等提示与视频历史联系起来，输出后续帧中的目标掩码。这里的历史记忆保存供当前分割读取的信息，不等于已经认定了现实车辆的永久身份；遮挡、离开和重新出现仍需分别检查\scite{vision-sam2}

长录像中的问题还可能回到早先片段。MovieChat于2023年提交预印本、2024年发表于CVPR，研究把密集视觉信息整理成长短期记忆以支持问答；压缩历史使较长输入可被利用，但压缩后的表示是否保留后来查询所需细节，仍取决于实际保存的内容。它扩展的是历史可用范围，不是给任意过去问题都提供完整证据\scite{vision-moviechat2024}

空间结果和响应形式也在扩展。2025年的视觉几何Transformer（Visual Geometry Grounded Transformer, VGGT）是联合几何预测模型，可从多视图形成相机、深度和点图等量，点图（point map）为图像位置提供对应三维点；Qwen2.5-Omni技术报告则把图像、音画输入与文字和语音响应联系起来。前者需要检查坐标、尺度和几何相容，后者还要检查答案内容、发音和响应时序，不能用一个共同接口分数概括这些不同任务\scite{vision-2503-11651,vision-qwenomni2025}

生成系统也从单幅图片扩展到视频、空间内容与音视频。共享模型和多阶段制作同时存在：一种路线共享部分表示与主干，另一种路线让专用模型通过明确接口传递结果。仅凭一个产品入口可以接收多种输入，无法判断内部是否共同训练、是否共享参数，更无法推出各项输出之间已保持一致。

\section{当前前沿与研究方向}
\label{sec:vision-overview-frontiers}

开放、连续和多源接口扩大了应用范围，也使证据保存、主体一致和成果协调成为更具体的要求。前沿问题可以按它们所缺的条件分开：有的缺观测证据，有的缺跨输出或跨模态的一致，有的缺对生成内容的控制，有的缺在有限预算下维持长期状态的能力。以下以截至2026年10月17日核对的代表文献定位这些问题，区分预印本首稿、修订和会议年份；这不是领域材料的穷尽。近期方法只用于说明改变了什么信息条件或状态组织，其局部结果仍限于对应数据、版本与观察预算。

\subsection{高分辨率与长时间观测的证据保存}
\label{sec:vision-overview-evidence-frontier}

整幅页面压缩到有限网格后，小字可能消失；长视频均匀采样后，短暂动作可能没有进入输入。研究问题是怎样先发现可能有用的时空范围，再读取足够细的原始内容，并保存以后可以回查的依据。网络中的视觉词元（visual token）是代表一个图像块或汇总范围的计算单元，增加这类单元的容量并不自动补回抽帧时删除的时刻。输入容量、证据选择和实际保留内容因此需要一起记录。图像块和时空块怎样形成视觉词元，在\ref{sec:vision-ivm-pipeline}节说明。

在前面的采样例子中，若查询是“那辆车的车牌是什么”，重新读取4.2至4.6秒的画面可以补充证据；但前提是原记录仍可访问，并有线索把查询引向这段时间。只保存“路口曾有车辆”的语言摘要，没有保存车牌区域，就不能从摘要恢复字符。主动再次观察与增加历史摘要解决的是不同缺口，也承担不同的读取、存储和等待成本。

VideoZeroBench的2026年v1为长视频问题标注答案、时间区间及空间框，并通过不同层级分别提供证据提示或要求模型返回证据。由此可以区分“给出证据后能否回答”和“能否自己找到证据并回答”，而不只看一个答案分数。其有限基准用于诊断这些环节，不证明全部长视频任务具有相同难度；证据获取和评价在视频感知章及\ref{sec:vision-understanding-evaluation}节展开\scite{vision-videozerobench2026}

\subsection{空间、时间与主体状态的一致性}
\label{sec:vision-overview-consistency-frontier}

相机、深度、点图与轨迹若描述同一场景，应在共同坐标及尺度约定下相容。一个输出把车辆放在路灯前方，另一个输出却给出相反的深度关系，即使各自看起来合理，组合结果也不能直接使用。VGGT提供联合几何预测的路线，但共同预测不等于所有输出自动相容；空间问答又检查模型能否依据观测回答方向、距离或数量问题。几何误差与答案错误需要分别评价，坐标和尺度约定也需要恢复到实际使用的场景。

时间上的一致性也有不同对象。车辆遮挡后再次出现，需要核验身份；动作反向播放，需要核验事件顺序；生成视频中的物体离开画面后重现，需要核验主体与状态保持。这些问题都涉及连续性，却不能由一个通用的“稳定分数”替代。

例如，轨迹每帧只移动少量像素，仍可能在两车交会时交换身份；生成车辆外观始终相同，仍可能在离开镜头期间停在旧位置。研究需要将位置变化、身份延续和状态演化各自对应到可检查的量，再检查它们能否共同描述一个过程。图像空间恢复、视频对应和生成历史分别承担其中的计算，不能由视觉平滑推出物理过程正确。

\subsection{不同模态间的事件对应}
\label{sec:vision-overview-alignment-frontier}

声音和画面可以互补，也可能冲突。画面中的人嘴唇在动，不表示录音一定来自此人；字幕出现一个名称，也不证明该对象可见。模型需要保留谁提供了什么信息、发生在何时，再决定如何融合。

多源模型提供了联合计算的接口，但评价还要检查删除、替换或错配某种模态后的变化。模态消融（modality ablation）是在约定设置中移除或改变一种输入，再比较结果，以诊断系统使用了哪些信息。若问题单凭字幕就能回答，原本的准确率不能独自证明视觉路径有效；若替换鸣笛声后答案改变，也只能说明系统对这项替换敏感，尚未说明它找对了声源。

替换录音可能同时改变音量、环境噪声和输入分布，删除字幕也可能造成训练中少见的输入格式。核验时应尽量固定问题、画面和其他条件，记录究竟改变了什么，再将结果变化与时空对应是否正确分开检查。跨模态接口由\ref{chap:vision-multimodal-models}章解释，来源核验和错配诊断由多模态理解章展开。

\subsection{生成内容的控制与协调}
\label{sec:vision-overview-generation-frontier}

“一辆红色汽车在蓝色路牌左侧”同时要求对象存在、数量正确、属性归属和空间关系成立。局部修改还要求未指定对象保持，视频增加持续身份和动作，声音增加声源与时间对应。即使整幅图与请求的语义相近，也可能把红色错赋给路牌；评价需要核查关系的两端，而不只计算整图与文字的相似程度。

Show-o2的2025年v3研究在共享主干中结合文字预测和视觉生成，以不同输出模块形成文字及图像或视频，具体能力取决于相应变体。它为“解释参考画面后按条件制作”提供了能力组合的路线，不能仅凭共享主干推断解释必然准确或全部修改条件都被遵守\scite{vision-showo22025}

LTX-2的2026年v1则以相互连接的音频和视频分支共同形成媒体，使两种输出可以在生成过程中交换信息。这不同于先固定完整视频再补声音，也不同于Show-o2的文字与视觉输出组合。对于“车辆转弯后鸣笛两次”，验收要同时检查车辆、动作区间、两次声音及其对应，不从自然画面或同步短片段直接推出长叙事、物理接触及全局状态都正确；模型接口由\ref{chap:vision-multimodal-models}章解释，条件遵循、未修改内容保持和跨成果协调由图像生成及多模态生成章展开\scite{vision-ltx22026}

\subsection{有限预算下的持续观察}
\label{sec:vision-overview-streaming-frontier}

车辆正在经过路口时，在线系统不能先看完未来的转弯再声称即时作出判断。沿用前文的在线输入条件，因果性（causality）在这里约束的是当前输出不能读取尚未到达的观测。为明确允许读取的范围，设第$t$帧的拍摄时间为$\tau_t$、到达系统的时间为$a_t$，两者以同一时钟的秒数表示。在当前时刻$\tau$可使用的观测集合为
\begin{equation}
  \mathcal{O}(\tau)
  =\{(\bm{X}_t,\tau_t,a_t):a_t\leq\tau\}\text{。}
  \label{eq:vision-causal-observations}
\end{equation}
集合条件按到达时间筛选，而不是按帧编号筛选；在忽略传输与采集延迟的理想情形下，才有$a_t=\tau_t$。这个集合说明哪些信息已经到达，不保证系统仍完整保存它们；实际输入还受所保存的历史内容和读取预算限制。

若算法必须等待$\delta$秒后才能形成输出，它在时刻$\tau+\delta$可以使用更晚到达的观测，但这段等待应计入响应延迟（response latency），即从请求发出到约定结果就绪所经历的时间。图\ref{fig:vision-observation-window}在忽略传输延迟的情形下区分这些范围。离线处理只读取已经给定的完整记录，不是能够取得任意现实未来；对于先前的查询时刻，它可以使用记录中更晚的帧。

\begin{figure}[htbp]
  \centering
  % 时间坐标只作范围示意，不表示某个实验的采样频率。
\begin{tikzpicture}[x=1cm,y=1cm,font=\figurefont\small]
  \fill[FigureInput!10] (0,0.15) rectangle (4.6,1.3);
  \fill[FigureModel!10] (4.6,0.15) rectangle (7,1.3);
  \fill[BookPaper] (7,0.15) rectangle (10.2,1.3);
  \draw[-{Stealth[length=2mm]},BookInk,line width=0.8pt] (0,0) -- (10.5,0)
    node[below left] {观测时间};
  \foreach \x in {0.4,1.4,2.4,3.4,4.4,5.4,6.4,7.4,8.4,9.4}{
    \fill[BookInk] (\x,0) circle (1.5pt);
  }
  \draw[BookInk,line width=0.9pt] (4.6,-0.25) -- (4.6,1.5);
  \draw[BookMuted,dashed,line width=0.8pt] (7,-0.25) -- (7,1.5);
  \node[below] at (4.6,-0.25) {$\tau$};
  \node[below] at (7,-0.25) {$\tau+\delta$};
  \node[align=center] at (2.3,0.75) {已到达历史\\即时可用};
  \node[align=center] at (5.8,0.75) {等待期间\\新增观测};
  \node[align=center] at (8.6,0.75) {其余未来\\离线才可用};
\end{tikzpicture}

  \caption{同一查询在不同信息条件下可读取的帧，示意中忽略传输延迟。圆点表示实际观测，竖线是查询发起时刻$\tau$；在线即时处理只读取左侧，允许等待$\delta$时才可加入中间区间，离线处理还可使用右侧已给定记录中的后续帧。颜色之外用文字和线型区分范围。}
  \label{fig:vision-observation-window}
\end{figure}

循环结构用先前内部状态处理当前输入，缓存（cache）保存供后续计算读取的已有信息，长记忆扩大可保留的历史；这些名称都描述计算组织，因果性仍取决于实际读取了哪些时刻。例如，用双向读取得到的特征更新一个循环状态，该状态仍含有未来信息，不能因外层具有循环结构就称整个过程为即时在线。

Flash-VStream的2025年路线分开处理持续到达的帧和用户问题，用概括历史的信息选择需要保留或取回的细节。TempoGround的2026年v1则面向流式语言定位，关联相邻帧对象，预测进入、持续或离开的存在状态，再形成二维框及相机坐标系中的三维框。前者关注历史怎样供问答使用，后者关注语言目标怎样随到达画面保持对应；两者都需要记录采样、历史容量、观察范围及响应时间，具体机制由视频感知和多模态理解章展开\scite{vision-flash-vstream2025,vision-tempoground2026}

响应预算还要明确结束点：第一文字单元出现、完整答案就绪和全部媒体播放完成不是同一时刻。存储更多摘要可能保存较长经历，却仍遗漏小目标；处理更多区域可能取得细节，却增加读取和等待。在线结果用已到达证据判断当前或过去，未来预测则估计尚未观测的状态，两者可以出现在同一系统中，但应分开标明和评价。

\subsection{可查询场景状态的维护}
\label{sec:vision-overview-queryable-state-frontier}

设街景中的一辆车离开镜头，几秒后又从另一个方向返回。“现在画面里有什么”“它上次在哪里”“回来的是不是同一辆车”和“离开期间可能去了哪里”需要不同信息。当前检测回答可见对象，历史记录保留过去观测，对象对应判断身份，离视期间的位置则只能由额外观测或预测约束。保存最后一幅画面不能自动说明这辆车此刻仍停在旧位置。

持续、可查询的场景状态（persistent, queryable scene state）指一种随观测更新、能按对象和时刻读取的内部记录或表示。它把身份、位置、属性、事件及其依据联系起来，使后续问题可以回到同一对象，并在新证据出现时修订原判断。这里的状态可以不完整，也可以保留多个可能解释；它不是要求系统精确存下现实世界的全部物理量。

为了区分“维护记录”和“回答当前问题”，可以用一个教学接口写出两步关系。设$\mathcal{S}_{t-1}$是处理第$t$帧前的状态，$o_t=(\bm{X}_t,\tau_t,a_t)$是已经到达的观测，$q$是查询，$\hat r_t$是返回的结果及证据索引，则
\begin{equation}
  \begin{aligned}
    \mathcal{S}_t&=U(\mathcal{S}_{t-1},o_t),\\
    \hat r_t&=R(\mathcal{S}_t,q)\text{。}
  \end{aligned}
  \label{eq:vision-overview-state-update-query}
\end{equation}
$U$表示根据新观测更新状态，$R$表示按查询读出结果；$\mathcal{S}_t$是组织多项信息的结构，不预设为固定维向量。初始状态可以为空，也可以含有明确来源的既有地图。这是问题分工的示意，不是已经实现的通用算法：每个具体系统还须规定更新哪些记录、如何处理冲突、删去什么，以及无法确定时怎样返回。

沿车辆例子复核这两步，首次看见车辆时可以建立临时身份，保存区域和时刻；离开镜头后，把可见状态改为“当前未见”，保留最后观测及可能运动的预测，但不把预测伪装成新观测。车辆返回时，再用新外观、位置和其他可用证据核查身份；若原先关联错误，应修订身份关系并保留修订依据。用户随后问“上次在何处”应读历史记录，问“现在在哪里”则需要当前证据或明确注明估计，两个查询不能都返回同一旧框。

V-JEPA 2.1的2026年v1提供了特征预测的相关路线：训练时从部分可见图像或视频预测完整输入的内部表示，并把预测约束扩展到可见位置和多个表示层，以支持局部空间结构与时间信息的学习。这里的潜在特征（latent feature）是网络内部对观测的数值表示，不是直接可检查的实体表。该路线为视频感知提供表示与预测依据，其结构与训练在\ref{sec:vision-ivm-vjepa}节展开；但不能仅由特征预测推断系统已能维护现实对象的永久身份或回答任意状态查询\scite{vision-vjepa21-2026}

InfiniteVGGT的2026年v1从另一端研究持续几何：在因果图像流中保留首帧参考信息，并按预算选择滚动保存的历史特征，持续形成相机、深度与空间结果。它改变的是几何计算如何在有限内部缓存下延续，不是无界保存全部历史；有限长序列上的检验也不构成无限时间无漂移的保证。其空间恢复机制归入图像识别与检索章，与视频感知中的实体历史相互补充\scite{vision-infinitevggt-2026}

LiveWorld的2026年v1则研究生成环境中的离视动态（out-of-sight dynamics），即对象不在当前视野内时仍可能继续发生的变化。它将静态背景与动态实体分开维护，为局部活动实体设置虚拟观察位置，生成离视期间的演化，再将更新状态用于返回视角的画面。这个过程模拟未观测内容，不是从真实录像恢复了确定发生的事件；动态区域数目、局部演化和几何条件也限制其范围。其机制由多模态生成章展开，不能与真实历史证据混用\scite{vision-liveworld-2026}

《Generalist Open-World Temporal Perception》的2026年v1是一篇观点论文，提出让感知状态具有持续性、可寻址性、跨输出相容性、来源对应、可修订性及不确定性表达。可寻址性（addressability）在这里指后续查询能够再次指向同一实体、表面或事件，而不是只能读取一段无法定位的摘要。论文将这些要求组织为研究议程，并讨论相关组件证据；它不是已实现一个满足全部要求的通用系统的实验报告\scite{vision-generalist-temporal-perception-2026}

这些路线共同指向如何组织可持续使用的信息，却不处理同一种状态。研究中的核验至少要分清真实观测是否可回查、身份是否可修订、几何是否在共同参考下相容，以及未观测演化是否明确作为预测或创作。保留内部缓存、答对一次问题或生成一段连贯视频，只分别说明其中一部分；要支持持续查询，还需在新证据、遮挡、重访和预算变化下检查这些关系是否仍成立。

\section{本章小结}
\label{sec:vision-overview-summary}

街景上的分类、定位、增强与编辑之所以需要不同解法，是因为它们要求不同的成果，也容许不同的内容变化。明确观测与输入依据、输出形式与结果约束之后，才能判断模型需要保留什么表示、取得什么监督，以及怎样恢复和核验结果。图像建立空间关系，视频增加身份、时间与信息到达条件，多模态处理再连接查询、不同来源及不同输出。

研究范围从人工设计的观测约束，经学习式输出组织和开放条件，扩展到连续、多源与交互应用，并没有取消证据边界。采样可能遗漏内容，预测可能依赖先验，生成可能补出没有观测支持的细节；持续场景状态还要把过去观测、当前判断和未观测预测分开保存，并允许新证据修订原关联。今后的证据保留、跨模态对应、生成控制和长期状态研究，需要同时考虑能回答什么、依据是否可查以及预算是否允许，而不只扩大一个输入接口。

本章对开篇街景请求的共同回答是：任务规定需要形成的成果，来源与预算规定可以使用的信息，核验检查成果是否仍指向正确对象、区域、时刻及用途。本章还没有回答这些成果由什么计算形成。接下来的五章先回到经典基础：\ref{chap:vision-image-formation}章说明图像数组怎样由光、传感器和相机几何形成，随后各章讲变换与滤波、特征与分割、视频的表示与运动，以及多模态数据的对齐与融合。第二部分的\ref{chap:vision-image-video-models}章再说明图像和视频怎样被学习模型编码、怎样解码成任务输出，\ref{chap:vision-multimodal-models}章连接文字与声音；后三个部分的任务章在这些模型之上追踪监督、参数来源和结果恢复。只有这些关系明确，分类、定位、复原、编辑与持续查询的能力声明才有可比较的含义。

\paragraph{扩展阅读}
想系统补齐本卷第一部分的经典基础，可以先读几部教材。Gonzalez与Woods的《数字图像处理》按灰度变换、滤波、复原、形态学和分割组织方法，与\ref{chap:vision-image-transforms}章的内容最接近；Szeliski的《计算机视觉：算法与应用》第二版以特征、运动、三维重建和识别为主线，兼顾经典算法与学习方法；Forsyth与Ponce的《计算机视觉：一种现代方法》更多从成像物理、几何与概率推断讲早期视觉；Hartley与Zisserman的《计算机视觉中的多视图几何》系统推导射影几何、基础矩阵和光束法平差，是\ref{chap:vision-image-formation}章相机几何部分的标准参考\scite{vision-gonzalez-woods2018,szeliski2022computer,vision-forsyth-ponce2012,vision-hartley-zisserman2004}

学习式路线方面，Prince的《计算机视觉：模型、学习与推断》用概率模型统一表述视觉问题，处在本章所说人工设计规则与端到端学习之间；Goodfellow等人的《深度学习》在卷积网络一章解释局部连接、参数共享和池化为何适合网格数据，可作为\ref{chap:vision-image-video-models}章的预备；Torralba、Isola与Freeman的《计算机视觉基础》从成像和信号处理一路讲到神经网络、表示学习与生成模型，用同一套讲法连接本卷经典基础与学习模型两部分\scite{vision-prince2012,goodfellow2016,vision-torralba-isola-freeman2024}

关于研究历史，Marr在1982年出版的《视觉》把视觉看作从图像到形状描述的分层计算，并区分计算理论、表示与算法、物理实现三个分析层次，本章先规定成果、再选择计算的思路与这种分析方式相近；Andreopoulos与Tsotsos回顾了约五十年的物体识别研究并讨论今后方向，可补充本章历史一节只列代表工作的不足\scite{vision-marr-vision1982,vision-andreopoulos-tsotsos2013}多模态方向，Zhang等人综述视觉语言模型的预训练目标、迁移方式和在分类、检测、分割上的评测，Yin等人综述多模态大语言模型的结构、训练数据与评测，分别对应本章所说CLIP式的匹配接口和生成式视觉语言模型\scite{vision-vlm-survey-zhang2024,vision-mllm-survey-yin2024}
